%%%%%%%%%%%%%%%%%%%%%%%%%%%%%%%%%%%%%%%%%%%
% Página de resumen del proyecto en inglés%
%%%%%%%%%%%%%%%%%%%%%%%%%%%%%%%%%%%%%%%%%%%
% !TEX root = A0.MiTFG.tex

\pagestyle{fancy}
\addstarredchapter{Abstract}

\begin{center}
	\scshape
	E.T.S. de Ingeniería de Telecomunicación, Universidad de Málaga
\end{center}

\bigskip

\begin{center}
	\Large \scshape
	\textbf{\tfgtitlenameENG}
\end{center}

\bigskip \bigskip \bigskip

\begin{minipage}{\textwidth}

\textbf{Autor:} \tfgauthorname

\medskip

\textbf{Tutor:} \tfgtutorname



\medskip


\textbf{Department:} Ingeniería de Comunicaciones

\medskip

\textbf{Degree:} Grado en Ingeniería Telemática

\medskip

\textbf{Keywords:} optical communications, ultrashort pulses, chromatic dispersion, atmospheric turbulence, oceanic turbulence, unguided media propagation, Matlab.

\bigskip \bigskip


\end{minipage}

\begin{center}
	\textbf{Abstract}
\end{center}


This project presents a theoretical and simulation-based study of the broadening and distortion effects experienced by ultrashort optical pulses during their propagation through atmospheric and oceanic media, caused by chromatic dispersion and turbulence. Optical communications have attracted considerable interest as they provide data rates far exceeding those of acoustic or radio-frequency technologies, although their implementation in unguided environments still poses significant challenges, particularly in the ocean.
Using Matlab simulations, the impact of chromatic dispersion, responsible for temporal pulse broadening and limitations in transmission capacity, and turbulence, introducing phase and amplitude fluctuations that reduce the collected power at the receiver, is analyzed. The results obtained provide a better understanding of these limitations and their implications for the design of optical communication systems in atmospheric and underwater environments. 


\blankpage
